%
% display_crash_loop_buggy.tex
%
% THE FIDELITY NEGATIVE CONTROL for display_crash_loop.tex.  This is the
% canonical spec with exactly one change: the quarantine clause is struck from
% CrashRender, so a crash tallies the death but never quarantines the scene.
% With that one clause gone the crash-respawn loop becomes reachable, and ProB
% reaches "crashes > THRESHOLD * card(poison)" with the exact loop trace.  A
% model that cannot reproduce the known defect proves nothing; this file is the
% proof that the canonical model detects the loop the quarantine rule prevents.
%
% Do NOT treat this as a design artifact -- it models the BUGGY design on
% purpose.  The canonical, verified model is display_crash_loop.tex.
%
% Type-check:  fuzz -t docs/display_crash_loop_buggy.tex
% Model-check (must report: goal FOUND -- the loop):
%   probcli docs/display_crash_loop_buggy.tex -model_check -nodead \
%       -goal "crashes>THRESHOLD*card(poison)" \
%       -p DEFAULT_SETSIZE 2 -p MAX_INITIALISATIONS 20 -p MAXINT 8
%

\documentclass[a4paper,10pt,fleqn]{article}
\usepackage[margin=1in]{geometry}
\usepackage{fuzz}
\usepackage[colorlinks=true,linkcolor=blue,citecolor=blue,urlcolor=blue]{hyperref}

\begin{document}

\title{The lux Display Crash-Loop Quarantine: Fidelity Negative Control}
\author{Buggy variant --- quarantine clause struck from CrashRender}
\date{July 2026}
\maketitle

% ============================================================
\section{Introduction}
% ============================================================

A scene that passes Hub-side self-validation can still crash the Display's
renderer.  When it does, the Display process dies --- but the scene does not
die with it.  It stays live in \texttt{HubDisplay}, the authoritative store, so
the Hub's crash recovery re-marks it for repaint, the respawned Display receives
it, and it crashes again.  Each respawn opens a fresh window that steals
keyboard focus.  Nothing in this cycle attributes the death to the scene that
caused it, so nothing stops the loop; the operator observed it live during the
B5 demo, once every ten seconds, until the scenes were cleared by hand.

The fix is a \emph{quarantine} rule: after a scene is attributed as the cause of
repeated Display deaths, the Hub stops replicating it while keeping it in the
store for inspection.  This document models the crash/attribute/quarantine
lifecycle as a finite state machine, so ProB can confirm the two properties that
matter --- a quarantined scene is never replicated, and the number of crashes is
bounded (the loop terminates) --- and can reproduce the unbounded loop when the
quarantine rule is removed.

This spec is a companion to \texttt{display\_lifecycle.tex}, not an extension of
it.  That spec proves that exactly one Display process owns the socket path
under a two-lock discipline; this spec takes that single-owner Display as given
and models what happens when the scene it renders is poisonous.  The two
concerns are orthogonal: the bind-race spec governs \emph{who owns the socket},
this spec governs \emph{whether a poison scene can loop the Display}.  Keeping
them separate leaves the ProB-verified bind-race model untouched.

% ============================================================
\section{Basic Types}
% ============================================================

\begin{zed}
  [SCENE]
\end{zed}

\texttt{SCENE} is the set of scenes an agent has installed in the store.  Two
scenes suffice to exhibit a poison scene alongside a healthy one; three exercise
attribution against several poison scenes at once.  The carrier is bounded by
ProB's \texttt{DEFAULT\_SETSIZE}.

% ============================================================
\section{Free Types}
% ============================================================

\begin{zed}
  DISPLAY ::= up | down
\end{zed}

The state of the single Display process: \texttt{up} while it is serving,
\texttt{down} after a renderer crash has killed it and before it is respawned.
That there is exactly one Display process is the guarantee proved in
\texttt{display\_lifecycle.tex}; this model represents that one process's
liveness and does not re-derive the singleton property.

% ============================================================
\section{Constants}
% ============================================================

\begin{axdef}
  THRESHOLD : \nat_1
\where
  THRESHOLD = 2
\end{axdef}

\texttt{THRESHOLD} is the number of attributed deaths a scene must cause before
it is quarantined.  It is two, not one: a single death admits a non-scene cause
(memory pressure, a driver fault) that happened to coincide with a render, and
must not quarantine an innocent scene.  Requiring the \emph{same} scene to be
the sole suspect across two deaths excludes the one-off, while a genuinely poison
scene --- which crashes the renderer on every render --- reaches the threshold
deterministically.

% ============================================================
\section{State}
% ============================================================

The state is flat: one schema, five components.  Its predicate carries only
\emph{structural} invariants --- a scene is quarantined only if it is poison and
only once its tally has reached the threshold.  These hold by construction in
every operation, so they act as coherence constraints, not as the safety
property under test.  The safety property (Section~\ref{sec:inv}), that the
crash count is bounded, is deliberately \emph{absent} here, so that a violating
state remains reachable and ProB can find it in the buggy variant.

\begin{schema}{Sys}
  disp : DISPLAY \\
  poison : \power SCENE \\
  quarantined : \power SCENE \\
  tally : SCENE \fun \nat \\
  crashes : \nat
\where
  quarantined \subseteq poison \\
  \forall s : SCENE @ s \in quarantined \implies tally~s \geq THRESHOLD
\end{schema}

\texttt{poison} is the set of scenes whose render crashes the Display --- an
environmental fact fixed for the run, never mutated by any operation.  A scene is
\texttt{poison} in the model exactly when the tree the agent installed hits the
renderer defect; a healthy scene is simply not in this set.  \texttt{quarantined}
is the set the Hub has stopped replicating.  \texttt{tally} counts the deaths
attributed to each scene, and \texttt{crashes} counts Display deaths in total ---
the quantity whose boundedness is the no-infinite-respawn property.

% ============================================================
\section{Initialization}
% ============================================================

\begin{schema}{Init}
  Sys'
\where
  disp' = up \\
  poison' \subseteq SCENE \\
  quarantined' = \emptyset \\
  tally' = (\lambda s : SCENE @ 0) \\
  crashes' = 0
\end{schema}

The Display starts up, nothing is quarantined, every tally is zero, and no crash
has happened.  \texttt{poison} is left free: ProB explores every valuation of
which scenes are poisonous, so the checks cover the healthy-only case, the
all-poison case, and every mix, in one run.

% ============================================================
\section{Operations}
% ============================================================

An agent renders a healthy scene without incident.  This operation changes
nothing but records, by its guard, that a scene outside \texttt{poison} never
crashes the Display and so is never quarantined:

\begin{schema}{RenderHealthy}
  \Xi Sys \\
  s? : SCENE
\where
  disp = up \\
  s? \notin poison
\end{schema}

The load-bearing operation is \texttt{CrashRender}: the Display renders a poison
scene that has not been quarantined, the renderer crashes, the death is
attributed to that one scene, and the scene is quarantined the instant its tally
reaches the threshold.  Folding the attribution and the quarantine into the
crash's own effect mirrors the design, where the recovery handler tallies and
quarantines synchronously before the next render --- there is no free-floating
``quarantine later'' step the scheduler could defer.

\begin{schema}{CrashRender}
  \Delta Sys \\
  s? : SCENE
\where
  disp = up \\
  s? \in poison \\
  s? \notin quarantined \\
  disp' = down \\
  crashes' = crashes + 1 \\
  tally' = tally \oplus \{ s? \mapsto tally~s? + 1 \} \\
  quarantined' = quarantined \\
  poison' = poison
\end{schema}

The guard \texttt{s? $\notin$ quarantined} is the whole quarantine defence: a
scene already quarantined can never be the subject of \texttt{CrashRender}, so it
can never crash the Display again.  The suspect here is a single scene,
reflecting the design's isolation mode --- after the first death the replicator
stops coalescing and sends one scene per cycle, so the scene in flight when the
Display dies is the sole suspect and the tally can only ever accrue against the
true culprit, never against an innocent scene that was co-batched with it.

A crashed Display is respawned.  This corresponds to the reap-then-ensure pair
that \texttt{display\_lifecycle.tex} proves yields exactly one socket owner; here
it simply returns the one Display to \texttt{up}:

\begin{schema}{Respawn}
  \Delta Sys
\where
  disp = down \\
  disp' = up \\
  poison' = poison \\
  quarantined' = quarantined \\
  tally' = tally \\
  crashes' = crashes
\end{schema}

The respawn backoff of the design --- which paces successive respawns so the
pre-quarantine deaths are not a rapid focus-stealing burst --- is a timing
policy, not a state-safety property: it changes \emph{when} \texttt{Respawn}
fires, never \emph{whether} a reachable state violates an invariant.  It is
therefore abstracted away here and specified in the design document, not modelled.

% ============================================================
\section{Invariants}
\label{sec:inv}
% ============================================================

Three properties must hold across all reachable states.  They are stated here as
predicates over \texttt{Sys}; how each is discharged is given in
Section~\ref{sec:verify}.  None is placed in the \texttt{Sys} schema: a predicate
placed there would be enforced as a guard on every successor, silently disabling
any operation that would break it and thereby masking the very defect this model
exists to catch.

\paragraph{Invariant 1 --- quarantined is never replicated.}
No quarantined scene is ever rendered.  This is a property of the transition
relation, established by construction: the only operation that sends a scene to
the Display and can crash it, \texttt{CrashRender}, is guarded by
\texttt{s? $\notin$ quarantined}, and \texttt{RenderHealthy} never touches a
poison scene at all.  A quarantined scene therefore never crosses to the Display.
ProB's operation-coverage report corroborates that \texttt{CrashRender} never
fires from a state in which its argument is quarantined.

\paragraph{Invariant 2 --- crashes are bounded (no infinite respawn).}
The total number of Display crashes never exceeds the threshold times the number
of poison scenes:
\[
  crashes \leq THRESHOLD * \# poison.
\]
Each poison scene can be the subject of \texttt{CrashRender} only while it is not
quarantined, and it is quarantined once its tally reaches \texttt{THRESHOLD} ---
that is, after \texttt{THRESHOLD} crashes attributed to it.  So each poison scene
contributes at most \texttt{THRESHOLD} crashes, and \texttt{crashes} increments
only in \texttt{CrashRender}.  This is the no-infinite-respawn property: the loop
can turn only a bounded number of times before every poison scene is quarantined
and the Display stays up.

\paragraph{Invariant 3 --- single Display.}
At most one Display process is live at any time.  This model represents the
Display as a single \texttt{DISPLAY}-valued component, so the property is
structural here; its real content --- that the reap/ensure respawn yields exactly
one socket owner and never two --- is proved in \texttt{display\_lifecycle.tex}
and inherited by composition.  \texttt{Respawn} is the abstraction of that
verified reap-then-ensure pair.

% ============================================================
\section{Verification}
\label{sec:verify}
% ============================================================

Type-checking is a \texttt{fuzz} pass:
\begin{verbatim}
  fuzz -t docs/display_crash_loop.tex
\end{verbatim}

Invariant 2 is a state property, checked by asking ProB whether its negation is
\emph{reachable}.  A goal that is not found means the invariant holds on every
reachable state:
\begin{verbatim}
  # Invariant 2 (must report: goal NOT found)
  probcli docs/display_crash_loop.tex -model_check -nodead \
    -goal "crashes>THRESHOLD*card(poison)" \
    -p DEFAULT_SETSIZE 2
\end{verbatim}

The model also reaches the stable fixed point in which every poison scene is
quarantined and the Display is up --- the state the loop settles into.  That this
state is reachable is a positive goal that ProB finds:
\begin{verbatim}
  # Stable state reachable (must report: goal FOUND)
  probcli docs/display_crash_loop.tex -model_check -nodead \
    -goal "disp=up & poison/={} & poison<:quarantined" \
    -p DEFAULT_SETSIZE 2
\end{verbatim}

Invariant 1 is a transition property, established by the guards of
\texttt{CrashRender} and \texttt{RenderHealthy} as argued in
Section~\ref{sec:inv} and corroborated by ProB's operation-coverage report.
Invariant 3 is inherited from \texttt{display\_lifecycle.tex}.

The \texttt{-nodead} flag is used deliberately.  Unlike the bind-race spec, whose
liveness depends on never deadlocking, this model is \emph{meant} to terminate:
the state in which every poison scene is quarantined and the Display is up has no
enabled operation, and that quiescence is the desired outcome, not a fault.  The
checks therefore search for the goal predicates rather than asserting
deadlock-freedom.

% ============================================================
\section{Fidelity: reproducing the crash-respawn loop}
\label{sec:fidelity}
% ============================================================

A model that cannot reproduce the known defect proves nothing.  The buggy variant
is the current specification with one change: the quarantine clause is struck from
\texttt{CrashRender}, so a crash tallies the death but never quarantines the
scene.  Concretely, replace the conditional assignment to \texttt{quarantined'}
with
\[
  quarantined' = quarantined,
\]
and, because the \texttt{Sys} coherence constraint
\texttt{quarantined $\subseteq$ poison} still holds trivially (the set stays
empty), the model still type-checks.  Now \texttt{quarantined} never grows, so the
guard \texttt{s? $\notin$ quarantined} never excludes a poison scene, and the
crash-respawn cycle

\begin{enumerate}
  \item \texttt{CrashRender}~$s$: the poison scene crashes the Display,
    \texttt{disp = down}, \texttt{crashes} increments.
  \item \texttt{Respawn}: \texttt{disp = up} again, \texttt{quarantined} still
    empty.
  \item \texttt{CrashRender}~$s$: the same scene crashes the same respawned
    Display.
\end{enumerate}

repeats without bound.  For any initial state with a poison scene, ProB reaches
\texttt{crashes > THRESHOLD * card(poison)} --- the same goal that is
\emph{not found} for the fixed design is \emph{found}, with the loop trace above,
for the buggy variant:
\begin{verbatim}
  # Buggy variant (must report: goal FOUND)
  probcli docs/display_crash_loop_buggy.tex -model_check -nodead \
    -goal "crashes>THRESHOLD*card(poison)" \
    -p DEFAULT_SETSIZE 2
\end{verbatim}

Equivalently, the stable-state goal
\texttt{disp=up \& poison/=\{\} \& poison<:quarantined} is \emph{not found} for the
buggy variant --- with \texttt{quarantined} pinned empty, a non-empty
\texttt{poison} can never be fully quarantined, so the loop never settles.  That
difference --- the fixed design terminates and the buggy design does not --- is
the evidence that the model detects the loop the quarantine rule was added to
prevent.

\end{document}
